\section{Introduction}


\subsection{Background and Motivation}

* Energy transition imperative: Climate goals driving decarbonization (breif)

* Renewable energy integration challenges: VRE variability creating stability and flexibility needs

* EV adoption trends and grid implications: Growth creating both challenges and opportunities

* EVs as distributed flexibility resources: technical potential for grid services

* Triple benefits of EV flexibility: Environmental, infrastructural and economic advantages

* NTNU campus case study relevance: Why this specific context matters (because its a separate microgrid?)

\hrule

"There is no denying the fact that electricity is crucial for the world economy to thrive in this and the coming centuries. The demand for electricity increases not only because the human population increases, but also because the social-economic activities of humans are rapidly shifting from manual processes to automated processes, which are essentially driven by electricity. Therefore, electricity becomes inevitable for the sustainability of modern civilization because it has found its way to the root of many human establishments [1,2]. Electricity can be generated from various sources and its generation from many sources is desired to boost the total generation capacity to meet the ever-growing demand. As electrical power cannot be stored, instantaneous balancing becomes a necessity and it is considered a key design parameter in power systems [3]." \cite{salman_review_2022}

"The evolution of the power system has significant impact on the operation and planning of the future power system. Five major global trends influencing the evolution of the power system are [1]: Decarbonisation - Decreasing the carbon footprint from electric power production. Decentralisation - Transition from few and large, centralized, power plants to many smaller, decentralised, power production units. Integration – Increasingly integrated electricity markets, greater interconnection of previously independent grids, and more integrated energy systems including sector coupling. These trends bring challenges in planning and operating the power system in a secure and reliable manner, such as: Identification of true operational state - As the utilization of the power system is increasing, it becomes even more important to quantify uncertainties in measurements and modelling to better understand reserves to critical limits. In particular to maintain the stability of the power system with regard to: frequency, voltage and rotor-angle stability. Changes in dynamic response - As a result of power electronic (PE) interfaced devices taking over the roles of rotating machines, the dynamic behaviour of the power system changes and, as a result, commonly accepted rules and principles may no longer be valid. New utilization patterns - More dispersed generation units, new type of demand, increased number and size of interconnections, results in utilisation of the power system in new and previously unforeseen ways." \cite{emil_hillberg_flexibility_2019}.

"Electrification is placing an additional demand on the power grid, but is also able to provide more demand-side flexibility with smart scheduling. EV's may also provide a storage capacity to decouple the production and demand timings. With limited supply-side flexibility (with the decrease in dispatchable coal and gas and increase in variable solar and wind), these demand-side flexibility options are becoming more and more important. (example: water heaters). In the future, electrification of heat and fuel may be able to provide long-term storage for dealing with seasonal imbalances (solar in Norway)" \cite{emanuele_taibi_power_2018}

"High VRE (variable renewable energy) increases the need to balance supply and demand. Flexibility is therefore increasingly important. "From power generation to stronger transmission and distribution systems, storage (electrical and thermal) and more flexible demand (demand-side management and sector coupling)." \cite{emanuele_taibi_power_2018}


"In [51], authors analyzed the impacts of DSF participation in a spot market and reserve market with respect to the German power network for the scenario of 2030, and a summary is provide in Table 4. It is observed that use of DSF decreases the impacts of reserve provision on the cost of a system, and the largest reduction is observed in the scenario of low coal/high VRE." \cite{kaushik_comprehensive_2022}

"V2G services could directly assist integration of renewable power. Kempton et al. (2005) calculated that V2G in the U.S. could stabilize large-scale (50\% of U.S. electricity) wind power with 3\% of the vehicle fleet dedicated to regulation for wind and 8–38\% to provide operating reserves or storage [315]. Ekman (2011) investigated the effects of different charging strategies on the balance between wind power production and consumption in a future Danish power system. Ekman finds that increased PEV penetration will reduce the excess of wind power, but PEVs alone will not be sufficient to fully utilize the wind power potential. Other balancing mechanisms will be needed if the wind penetration exceeds 50\% [307]." \cite{lund_review_2015}

"In domestic applications, Yoshimi et al. (2012) find that V2G can be used to increase the rate of utilization of PV-generated electricity, cut CO2 emissions and lower the costs of purchased electricity [317]." \cite{lund_review_2015}

"The societal advantages of developing V2G include an additional revenue stream for cleaner vehicles, increased stability and reliability of the electric grid, lower electric system costs, and eventually, inexpensive storage and backup for renewable electricity." \cite{kempton_vehicle--grid_2005}

"Furthermore, implementing V2G solutions with EV fleets can bring significant net value (i.e., negative incremental cost) to the system due to its ability to mitigate the reinforcement cost of electricity infrastructure and improve the economics and carbon footprint of meeting frequency regulation requirements. At the same time, V2G can significantly reduce carbon emissions from the power system, especially when this reduction is quantified per kilometre of distance driven, and considering that these benefits would be additional to other, more localized benefits such as improved local air quality." \cite{aunedi_whole-system_2020}


\subsection{Problem Statement}

* Dual challenge of energy transition: VRE variability and rising electricity demand

* Underutilized flexibility potential: Brief mention of technical vs. market gap (simplified, with a ref to the theory section which covers this)

\hrule
"The energy transition leads to a shift from conventional, controllable power plants to renewable, variable generation resources such as wind and solar power plants. According to [1,2], the existing renewable power capacity worldwide (including hydro power) nearly tripled from 1150 GW to 3146 GW between 2008 and 2021. The variability and uncertainty of renewable generation leads to an increased need for flexibility resources." \cite{lechl_review_2023}

"The increasing penetration of renewable energy sources (RES), particularly wind and solar power, has introduced significant variability into the power system. Unlike conventional generators, RES are subject to external factors such as weather conditions, resulting in fluctuating and less predictable power output. This variability complicates the balancing of supply and demand, as renewable generation often does not align with peak consumption periods or established demand patterns, thereby challenging traditional methods of maintaining grid stability" \cite{impram_challenges_2020}.

"Nordic power demand could double by 2050, with solar and wind projected to become the largest sources of power production." \& Demand and consumption flexibility must increase - EV's can contribute to both. \cite{noauthor_nordic_2023}


\subsection{Objectives/Research Questions}

* Primary objective: Evaluate EV flexibility potential for increasing campus grid performance

* Specific research question suggestions: \\
** Technical flexibility quantification under different EV penetration scenarios \\
** Grid stability and reliability impacts of V2G integration \\
** Economic viability and market accessibility gap assessment \\


\subsection{Scope and Limitations}

* Study boundaries: NTNU Gløshaugen campus focus, employee private vehicle and NTNU service vehicle charging.

* Methodological limitations: Modeling assumtions when creating the model in PowerFactory \& (no real-time market integration in the simulations)

* Data constraints: Historical charging data availability, generalization limitations, removal of edge-cases

\subsection{Thesis Structure}

* 1-2 sentences describing each chapter's contribution


\section{Theory}

\subsection{Power Systems, Renewable Energy Integration and Flexibility}


* Intro text (without a subheading of "traditional power systems":\\
** Historical grid paradigm: Centralized, reliable, dispatchable generation with controllable coal, natural gas, hydropower plants \\ 
** Traditional balancing and grid stability mechanisms: Synchronous turbines providing frequency stability and predictable dispatch \\
*** Automatic frequency response, physical inertia from rotating masses that resist frequency changes, reactive power support \& are grid-forming: Establish voltage and frequency through their rotating electromagnetic fields \\
** Transition challenge: Short sentence on how changing over to VRE challenges this previous stability paradigm

\hrule

"The power grid has essentially no storage (other than its 2.2\% capacity in pumped storage), so generation and transmission must be continuously managed to match fluctuating customer load. This is now accomplished primarily by turning large generators on and off, or ramping them up and down, some on a minute-by-minute basis." \cite{kempton_vehicle--grid_2005}

"The non-storability of electrical energy has dictated balancing as key design parameter of power systems. Flexibility is therefore an inherent feature of the planning and operation of power systems, which are designed to ensure spatial and temporal balancing of electricity generation and consumption at all times" \cite{g_papaefthymiou_power_2018}



\subsubsection{Variable Renewable Energy Sources: Proliferation and Challenges}

* VRE growth trajectory: Brief information on global adoption trends and penetration levels, and more focused on the Nordics adoption \& penetration. Focus on wind and solar (as they are growing the most?) Provide graphs and future predictions. Also emphasize the Nordic grid characteristics (high hydro)

* Technical grid challenges: \\
** Variability and forecasting uncertainty \\
** Reduced system inertia and stability services: Solar and wind is mostly grid-following technologies, and with increasing VRE penetration, the system looses important grid-forming capabilities. \\
** Reactive power and voltage support requirements. \\

* Operational challenges: \\
** Increased flexibility requirements ($>30\%$ penetration threshold \cite{alireza_akrami_power_2019}) \\ \
** Need for faster ramping and balancing services \\
** Grid congestion, curtailment, and security of supply issues from over-/under-production of VRE's

\research{Renewable peaking technologies vs. fossil peaking plants. For renewables: Are there are only hydropower, battery storage and (in the future - hydrogen}

\hrule

"The major share of RE is supplied from the wind energy and solar energy sources. Generation from these energy sources is variable and uncertain in nature. Hence, these sources are considered as variable renewable energy (VRE) sources. More than 50\% of Denmark’s electricity was supplied by wind energy and solar energy in the year 2020. Similarly, more than 30\% of Ireland’s electricity was supplied by VRE sources in the year 2020. China is planning to achieve VRE share of more than 60\% by 2050 and United States planned to achieve VRE share of 80\% by 2050 [2–4]. Hence, due to high share of VRE sources in future power generation scenarios, the high variability on the supply side has become a major concern" \cite{kaushik_comprehensive_2022}.


"Reference [116] has declared that the need for flexibility highly increases at a renewable penetration greater than 30\%, especially for solar generation. Also, a study of the Philippine Department of Energy estimates that about 50\% of extra generation capacity required as future reserve capacity must be from mid-range and peaking generation units [117]. An analysis in [118] shows that for 100\% renewable electricity supply scenarios at 2030 and 2050, a generation mix with a nominal capacity of more than twice maximum demand is required. All these studies reveal that more penetration is in parallel with the need for more flexibility." \cite{alireza_akrami_power_2019} 


"A critical review of the papers discussed in this section is carried out, and a comparative study is prepared for all types of drivers considering various aspects and dimensions which affect the PSF and are illustrated in Table 3. Results are expressed on a scale of 100. Higher values indicate the increased need for flexibility. Increased Solar PV generation at DSO level: 100. Increased wind generation at TSO level: 95. PV at TSO level: 87. Increased wind at DSO level: 86." \cite{kaushik_comprehensive_2022}.


"Physically, power network provides sufficient capacity for transferring power from generation plants to consumers, which is traditionally considered as a fixed structure. However, due to the integration of large-scale renewable generations, the operation of power network has been pushed much closer to its technical limits than before. As a consequence, new problems arise which can hardly be resolved when there is lack of grid-side flexibility. First, the actual outputs of uncertain renewable generation may remarkably deviate from the forecast, causing high operational risk or congestion. The congestion can further increase operating cost and limit the usage of available flexibility resource from the generation and/or demand sides. Second, renewable generation usually requires the support of sufficient reactive power, imposing additional risk of voltage instability on power system operation" \cite{li_grid-side_2018}

"Reduced system inertia, increased variability and uncertainty in generation, bi-directional power flows in distribution networks." \cite{emil_hillberg_flexibility_2019}

"A significant barrier to the integration of renewable generation has historically been the perceived need for additional capacity for periods when renewable generators are unable to produce. While an understanding of the capacity adequacy of variable generation (VG) is maturing [5], [6], system operators are now also concerned with the magnitude and uncertainty of the ramps associated with VG [7]-[10]." \cite{eamonn_lannoye_assessing_2015}

"electricity services [6], such as frequency regulation (FR). FR will be particularly critical in systems with very high penetrations of variable renewable sources such as wind and solar PV." \cite{aunedi_whole-system_2020}

"On long term, increasing renewable power will decrease the demand for base load power plants and increase the demand for peak load capacity [365]. This shift of conventional generation capacity and increasing fluctuation of net load will increase the price volatility [365]. As the predictability of variable renewable energy sources is low on the day-ahead market, the demand for balancing services and intraday trading will consequently increase [367,368]. The costs of RE and wholesale prices will also influence investment, expansion and retiring decisions of transmission and generation in the long run, making the average effect of RE on wholesale prices less clear than on short term [2]." \cite{lund_review_2015}


\subsubsection{Flexibility Definitions and Metrics}

* Definitional landscape and consensus of what flexibility is: \\
** Overview of varying flexibility definitions in literature \\
** Identification of core common elements across definitions \\
** Selection of working definition for this thesis analysis

* Components and categories of power system flexibility: \\
** Types of flexibility needs (power, energy, transfer capacity, voltage) \\
** Hierarchy from potential to market-accessible flexibility 

* Quantitative metrics for flexibility assessment: \\
** Technical dimensions (power, energy, ramp-rate) \\
** Operational characteristics (response time, duration, recovery) \\
** Directional considerations (up/down[supply/absorb] regulation capabilities

\research{When this section is a little bit more refined, check article 10 \cite{kaushik_comprehensive_2022}, which has a more meta-overview of the flexibility landscape}


\hrule

"The ability of power system operation, power system assets, loads, energy storage assets and generators, to change or modify their routine operation for a limited duration, and responding to external service request signals, without inducing unplanned disruptions" \cite{degefa_comprehensive_2021}

“Power system flexibility is defined as the ability of a power system to reliably and cost-effectively manage the variability and uncertainty of demand and supply across all relevant timescales”
\cite{international_energy_agency_status_2017}

"Definitions of flexibility that refer to the overall energy (or power) system usually describe it as the capability of that system to absorb disturbances sufficiently fast to maintain stability [7,13,14]." \cite{lechl_review_2023}

"The direction of the changes in net load is an additional consideration for flexibility. The magnitude and frequency of occurrence of net load changes, and resources available to meet upward and downward changes are asymmetric. For example, resources at maximum output can only assist when net load is decreasing, while offline resources may be able to come online during periods of increasing net load." \cite{eamonn_lannoye_evaluation_2012}

"A power system can be considered flexible if it can cost-effectively, reliably and across all time scales:
\begin{enumerate}
    \item Meet the peak loads and peak net loads, avoiding loss of load.
    \item Maintain the balance of supply and demand at all times, and ensure the availability of sufficient capability to ramp up and down, the availability of sufficient fast-starting capacity and the capability to operate during low net loads.
    \item Have sufficient storage capacity (both electricity storage and, through sector coupling, renewable heat and gas) to balance periods of high VRE generation and periods of high demand but low VRE generation.
    \item Incorporate capabilities to adjust demand to respond to periods of supply shortages or over-generation
    \item Maintain capabilities to mitigate possible events that could de-stabilize the power system through maintaining an adequate supply of ancillary services at all times.
    \item Operate under a well-designed market where existing flexibility is not locked by market inefficiencies (see section 3.2).
\end{enumerate}
" \cite{emanuele_taibi_power_2018}

\nb{Make sure that the points covering temporal and market aspects of flexibility get a /ref to the sections which will expand upon those points [this last definition should point to technical vs. market flexibility]}




"Flexibility needs are divided into four categories; flexibility for power, energy, transfer capacity, and voltage [10]. • Flexibility for power: The power supply-demand balance is needed to maintain frequency stability for short-term periods (a second to an hour). It is due to an intermittent weather condition-dependent power supply in generation. • Flexibility for energy: The energy supply-demand balance is needed for demand scenarios over medium to long term periods(hours to several years). It is due to a decrease in fuel storage-based energy supply in generation. • Flexibility for transfer capacity: Power transferability is needed to prevent bottlenecks over short to medium term periods (minutes to several hours). It is due to increased peak demands, increased peak supply, and increased usage levels. • Flexibility for voltage: The bus voltages need to be kept within predefined limits over short term periods (seconds to tens of minutes). It is due to increased distributed generation in distribution system resulting in bi-directional power flow and operation scenarios variance." \cite{impram_challenges_2020}

\nb{Make sure that the points covering temporal and market aspects of flexibility get a /ref to the sections which will expand upon those points [These quotes should point to flexibility requirements at different time scales]}

"In Ulbig and Andersson’s study [23], the sources of power system flexibility are divided into four categories: • Potential flexibility sources; are physically available and usable flexibility sources. However, they are not controllable or observable. • Actual sources of flexibility; are the usable part of potential flexibility sources because they are controllable and observable. • Flexibility reserves; are the economically usable part of the actual flexibility sources. • Flexibility reserves in the power market; are parts of the flexibility reserves that can be obtained from the power or ancillary services market." \cite{impram_challenges_2020}

"Resolving this ambiguity is a main objective of this work. To achieve this, a taxonomy called FLEXBLOX ("a box of flexibility building blocks") is introduced that clearly separates the terms flexibility resource, flexibility potential and flexibility requirements and mathematically defines their relationships, resulting in an unambiguous flexibility definition." \cite{lechl_review_2023} \\
* Definition 2: Flexibility metric for requirements and potentials have three measurements: apparent power $\sigma$ (in VA or W), energy $\epsilon$ (in Wh), and ramp-rate $\rho$ (in W/min) \\
* Definition 3: Flexibility requirement \\
* Definition 4: A flexibility resource \\
* Definition 5: A flexibility potential \\
* Definition 6: A flexibility mediator - grid or markets \\
* Definition 7: Flexibility - based on flexibility requirements by using flex resources' flexibility potentials, which are accessed by flexibility mediators \cite{lechl_review_2023} \\

\research{Read more about the FLEXBLOX taxonomy from article 3: \cite{lechl_review_2023} and expand upon the definitional concepts}
\illu{Figure 1 and 2 from Article 3 - \cite{lechl_review_2023} (specially fig. 2, provide a fantastic overview of the interconnectivity across flexibility categories. Use/make my own/write in text)}





"Operational flexibility is defined in terms of power capacity (MW), ramp rate (MW/min), i.e., the ability to increase energy production with a certain rate, and ramp duration (min), i.e., the ability to sustain ramping for a given duration" \cite{kaushik_comprehensive_2022}.

"Flexibility trinity: Power ramping capability (MW/min), Up- and Down regulation (MW), Storage capability(capacity, MWh)." \cite{kaushik_comprehensive_2022}

"The metrics for defining flexibility can be derived from these effects. Huber et al. (2014) [27] used three metrics to characterize flexibility requirements, namely ramp magnitude, ramp frequency and response time, in particular of the net load which results when the variable renewable generation has been subtracted from the gross load." \cite{lund_review_2015}.

From "2.5 Quantifying the need for flexibility in the power system" in \cite{emil_hillberg_flexibility_2019} \& \cite{degefa_comprehensive_2021} has basically the same definitions in their "Table 3".
* Quantifiable dimentions: \\
* Power (capability to deliver. Size of the active or reactive load (MW or Mvar) \\
* Response time/reaction time (start up time of a resource) (in seconds, minutes etc..) \\
* Ramp rate (MW/kW per second) \\
* Energy (capacity): time duration of a certain amount of power (kWh etc.) \\
* Recovery period: How long it takes to recover the energy used to provide flexibility. \\
* Service duration: How long the flexibility can be provided before the resource is spent. \\
* Ramp duration: Time needed from activation to ramping up to full capacity (seconds; The power capacity divided by the ramping capacity)
\cite{emil_hillberg_flexibility_2019}, \cite{degefa_comprehensive_2021} \\

\illu{From: Article 22, Figure 2 (page 6) - \cite{degefa_comprehensive_2021}. A good visualization of the most important flexibility/battery metrics}


\subsubsection{Sources of Flexibility}

* Comprehensive overview of flexibility source categories: \\
** Supply-side, demand-side, energy storage, grid infrastructure enhancements (interconnectivity of regions) \\
** Establish the four-fold texonomy from power system perspective (???) \\
** Keeping the discussion around flexibility-sources other than battery storage technologies fairly brief

* Definitional precision and conseptual boundaries: \\
** Clarify flexibility vs. DSM/DSR terminology confusion \\
** Distinguish true flexibility from energy efficiency and load shedding etc (\cite{degefa_comprehensive_2021}) \\
** Establish that load shifting and peak management constitute legitimate flexibility (especially for EVs) \\

* Energy storage systems as primary flexibility resource: \\
** VRE integration necessity and storage hierarchy across system levels \\
** Versatility advantages: Availability, response times, and capacity scalability \\

* Battery technology characteristics and economic constraints: \\
** Battery cost, life-cycles, performance (maybe compare different battery technologies) \\
** Economic limitations - capital costs vs. operational benefits \\
* Distribution-level applications (??)

* Technical grid service capabilities: \\
** Active and reactive power injection/absorption capabilities \\
** Power quality services: Voltage regulation, harmonics mitigation \\
** Integration benefits and necessity for renewable energy systems


\hrule


"Flexibility can be increased by demand side management (DSM), energy storage systems (ESS), energy curtailment, sector coupling, and expansion of transmission grid [26]." \cite{kaushik_comprehensive_2022}.

"Flexibility of power is provided by demand-side response (DSR), short-term battery energy storage system (BESS), virtual inertia, fast frequency response, and power system stabilizer. Flexibility of voltage is provided by onload tap changers of transformers, voltage boosters, automatic voltage regulators, and coordinated voltage control. Flexibility of transfer capacity is provided by the DSR, BESS local support, topology changes, rescheduling, and phase shifting transformer (PST). Flexibility of energy is provided by back-up generation, energy storage for long term support, optimization techniques, and HVDC supergrid [113–116]." \cite{kaushik_comprehensive_2022}

"There are terms which are often confused with flexibility, such as demand side response (DSR) / demand response (DR), demand side management (DSM), flexible generation and energy storage, on both the supply and demand side. These terms represent only parts of the definitions of flexibility and are not alternative terms. In [8], DSM is described as activities to activate the demand side, comprising actions such as energy efficiency, savings, self-production and load management. Further, load management techniques and DR are examples of DSM solutions. According to [38], there are six typical versions of DSR, such as conversion and energy efficiency, load shifting, peak clipping, valley filling, flexible load shape and electrification. The definition proposed in this paper assumes the following: * Energy efficiency is not flexibility [...]. * Load shedding is not flexibility [...]. * Curtailment of generation based of VRES is not flexibility [...]." \cite{degefa_comprehensive_2021}
\research{Read this section from article 22 \cite{degefa_comprehensive_2021} in detail - acknowledge that some parts of DSM (energy savings etc.) are not to be considered flexibility solutions - but that peak-shaving, valley filling and load shifting can shifting can still be flexibility solutions for EVs since it can be done without affecting the operability of the car. The main point of this article as I understood it was that measures which in any way affect the operability of buildings/industry etc. are not flexible. While an EV can provide this and still be at full charge by the end of the workday}

\research{Find a table or an illustration of some kind to clearly and succinctly be able to show different kinds of technology and which kinds of flexibility these technologies checks of on. [first column: technologies - rest of the columns: flexibility needs, and an X (or a colour scale rating) of how well the technology can support these]}

"Deployment of energy storage in parallel with the high penetration of VERs is undoubtedly needed. While the output power of VERs is accompanied with variability, energy storage can assist the power system to absorb the surplus generation of VERs in the case of over-generation or discharge of their energy to the system to help with any case of production scarcity. Further, storage might be performed in the ternary levels of the power system, such as water in pumped-hydro storage as a primary input at the generation, bulk battery storage at the transmission, and electric vehicles at the distribution levels [44]." \cite{alireza_akrami_power_2019}

"Energy storage technologies are also used in the power system to improve flexibility at large scales. However, use of these technologies is limited due to economy and technological constraints [57]." \cite{kaushik_comprehensive_2022} 

"Small-scale storage installed in distribution grids is an interesting flexibility option on the local level. Often fast and accurate, e.g. batteries, it is an ideal option to provide ancillary services. Although other flexibility options such as demand side flexibility are considered to be cheaper, there currently is a trend towards installing batteries, e.g. in the US and Germany" \cite{g_papaefthymiou_power_2018} 

"The lithium-ion battery is deployed widely in the market for small appliances. They have high efficiency, good energy density and low self-discharge rate. For the moment, though, they are still expensive for large-scale power [141,225]. Batteries can be considered as a kind of bridging technology between bulk storage and high-fidelity storage, e.g. for applications in which flexibility is desirable. Unfortunately, batteries are currently limited by their short operational lifetimes." \cite{lund_review_2015} 

"When energy storage units (ESUs) within the distribution grid, such as batteries, provide local services such as supporting the integration of photovoltaics, peak shaving, and infrastructure upgrade deferral, they are inactive or only partially used most of time. Moreover, they are often not profitable because of their high investment costs." \cite{olivier_megel_scheduling_2015}

"Energy storage units are considered the most versatile flexibility resource due to their high availability, fast reaction times and variable capacity sizes [52–54]. \cite{lechl_review_2023}

"Battery energy storage can both inject and absorb active and reactive power in the network to help solve under-voltage, over-voltage, voltage unbalance, power factor correction, harmonics and mitigate flicker." \cite{degefa_comprehensive_2021}

\illu{Take inspiration from illustration about "Value-stacking" from Article 21. \cite{noauthor_veileder_nodate} - Make/write a version of it which illustrates all the different kinds of flexibility battery storage offers}

"As most batteries have self-discharge losses as well, they are mainly feasible for short-term storage. Another disadvantage is that their performance reduces with increasing number of cycles. The capital cost and replacement cost of batteries dominates the cost of stored energy with batteries, while operation and maintenance costs are much less significant [217]. Batteries have near-instantaneous response times, which is a valuable feature for improving network stability [146] e.g. with renewables. The main use is providing power quality, short-term fluctuation reduction and some ancillary services or transmission deferral [22,78]. Batteries are modular in size enabling flexible siting e.g. close to load or production, or even changing location over the life-time [22]. Batteries for integrating variable renewable generation are already in use. For example in Futumata, Japan, a 51 MW wind farm is supported by a 34 MW sodium-sulphur-battery [218]. The benefits of a battery for PV and wind have also been verified through simulation studies [219–221]." \cite{lund_review_2015}

\nb{This above quote can be used as a preamble for the "electric vehicles as flexibility resources" section, by emphasizing that "the capital cost dominates the cost of stored energy with batteries" -- but all electric cars have already bought the car, and the battery is just sitting there being unused for most of the day}


\subsubsection{Flexibility Requirements at Different Time Scales}

* Introduction to temporal flexibility classification: \\
** Short-term - seconds to minutes, frequency regulation and immediate balancing (and more?) \\
** vs. medium-term - Minutes to hours, needed for managing intra-day variations (and more?) \\
** vs. long-term - Hours to days, important for seasonal variations and prolonged periods of high or low VRE output. \\
** vs. expansion planning(?), which is way outside the scope of this study

* Power quality and regulation services (milliseconds to 5 minutes): \\
** Frequency and voltage stabilization requirements \\
** High cycling capability needs vs. minimal energy storage requirements \\
** EV/battery suitability: excellent for power capacity, challenging for continuous availability

* Reserve services (5 minutes to 1 hour): \\
** Spinning vs. non-spinning reserve characteristics and requirements \\
** Capacity payment structures and economic advantages for EVs/batteries \\
** Regulation up/down market segmentation and automatic control requirements 

* Load management services (1 hour to 3 days): \\
** Load following, load leveling, and curtailment prevention functions \\
** Energy capacity requirements and duration constraints \\
** EV limitations for extended duration services

* Creating a comprehensive flexibility services table: \\
** Columns - service type, service category, response time, duration, power/energy needs, suitable technologies, EV/Battery suitability assessment \\
** Include something about market mechanisms (which market the different service requirements are traded in [spot, intra-day, balacing]), and compensation structures



"Having established the some of the sources capable of providing power system flexibility, particularly the characteristics of battery storage systems, we now examine how these resources map to specific operational flexibility requirements across different temporal scales."
\hrule

"The flexibility studies have been classified into two: short-term operational flexibility and long-term planning flexibility. A flexible power system should possess certain facilities that enable it to overcome uncertainties in addition to being able to compensate for deviations in generations and demands whilst maintaining the system’s reliability [11,12]." \cite{salman_review_2022}

"Power system flexibility analysis is needed both to ensure instantaneous stability of the power system and long-term security of supply. Additionally, the variability associated with VG can be separated into different timescales relevant to power systems, from seasonal and interannual variability to daily, hourly, and second-by-second variability (Mills, 2009)." \cite{michael_emmanuel_review_2020}




\emph{Power quality and regulation (millisecons to 5 minutes)}

"Power quality and regulation. Power intensive. Rapid and frequent response and very short duration. Used to balance fluctuations in network frequency and voltage that arise from variations in wind and solar generators' output." \\
"Energy storage can be used to mitigate these effects [141]. The storage systems are best suited for this service due to a rapid response time and high power ramping rate, as the fluctuations require action within seconds to minutes, and a high cycling capability, because continuous operation is required. A large storage capacity is here unnecessary as over 80\% of the power line disturbances last for less than a second [142]." \\
"a NERC report which claims that storage may not be a good replacement for the traditional stability services (system inertial response, automatic equipment and control systems), unless it also provides other grid services [22,78]." \cite{lund_review_2015} 

"Regulation is controlled automatically, by a direct connection from the grid operator (thus the synonym 'automatic generation control'). Compared to spinning reserves, it is called far more often (say 400 times per day), requires faster response (less than a minute), and is required to continue running for shorter durations (typically a few minutes at a time)." \cite{kempton_vehicle--grid_2005}


\emph{Spinning, non-spinning and contingency reserves (5 minutes to 1 hour)}

"Some markets split regulation into two elements: one for the ability to increase power generation from a baseline level, and the other to decrease from a baseline. These are commonly referred to as “regulation up” and “regulation down”, respectively. For example, if load exceeds generation, volt age and frequency drop, indicating that “regulation up” is needed. A generator can contract to provide either regulation up, or regulation down, or both over the same contract period, since the two will never be requested at the same time. Regulation is controlled automatically, by a direct connection from the grid operator (thus the synonym “automatic generation control”). Compared to spinning reserves, it is called far more often (say 400 times per day), requires faster response (less than a minute), and is required to continue running for shorter durations (typically a few minutes at a time)." \cite{kempton_vehicle--grid_2005}

"Spinning reserve refers to online power generation capacity synchronized to the grid having a short response time for ramping up but enabling several hours of use. They are generally used in contingency situations such as major generation and transmission failures [136]. Spinning reserves are restored to their pre-contingency status using replacement production reserves that should be online 30–60 minutes after the failure [136]. Non-spinning reserve is similar to spinning reserve, but without immediate response requirement. However, these reserves still need to fully respond within 10 minutes [78]." \\
"Performing a spinning reserve service requires both a rapid response time and a large capacity for storing energy. According to Rabiee et al. (2013), suitable technologies are batteries and flow batteries. As loads can shed very quickly, DSM is well-suited for reserve provision. Shiftable loads are in particular very suitable as the duration of the reserve provision is often short enough so that the load process is not disrupted [34]." \cite{lund_review_2015}

"Spinning reserves refers to additional generating capacity that can provide power quickly, say within 10min, upon request from the grid operator. Generators providing spin ning reserves run at low or partial speed and thus are already synchronized to the grid. (Spinning reserves are the fastest response, and thus most valuable, type of operating reserves; operating reserves are “extra generation available to serve load in case there is an unplanned event such as loss of generation” [16].) Spinning reserves are paid for by the amount of time they are available and ready. For example, a 1MW generator kept “spinning” and ready during a 24-h period would be sold as 1MW-day, even though no energy was actually produced. If the spinning reserve is called, the generator is paid an additional amount for the energy that is actually delivered (e.g., based on the market-clearing price of electricity at that time). The capacity of power available for 1h has the unit MW-h (meaning 1MW of capacity is available for 1h) and should not be confused with MWh, an energy unit that means 1MW is flowing for 1h. These contract arrangements are favourable for EDVs, since they are paid as “spinning” for many hours, just for being plugged in, while they incur relatively short periods of generating power." \cite{kempton_vehicle--grid_2005}




\emph{Load following, load leveling, curtailment prevention (1 hour to 3 days)}

"Load following is a continuous grid service that is used to obtain a better match between power supply and demand. Energy storage can be used for this purpose, by storing power during a period of low demand and injecting it back into the grid during a period of low supply [141]. Batteries and flow batteries, hydrogen, CAES and PHES are well-suited for this application [78,144]." \\
"Load leveling with energy storage refers to the evening-out of the typical mountain and valley-shape of electricity demand. As with load following, energy is absorbed during periods of low demand and injected back to the grid during high demand [141]. This allows baseload power generators to operate at higher efficiencies and also reduces the need for peaking power plants. This ancillary service can be provided by flow batteries, CAES, hydrogen and PHES [78,144], as all of them can handle large amounts of energy. Rabiee et al. also include batteries in this list [144]." \\
"Transmission curtailment prevention and transmission loss reduction are ancillary services that temporarily reduce the amount of current flowing in certain parts of the power grid, increasing the efficiency of transmission and preventing production curtailment due to power line limitations.With much renewable energy production and no means of storing excess power, power production may need to be curtailed (cut off) to ensure system stability or due to limitations in transmission infrastructure. However, with energy storage, the power plants may continue harvesting energy even while being disconnected from the rest of the grid. Renewable power is injected into the energy storage system instead of the grid, and when the grid is ready for the dispatch, the storage is discharged. The duration requirement for such measures ranges from 5 to 12 hours. Storage also allows increasing the efficiency of transmission. Because transmission losses are proportional to the square of the current flow, the net resistive losses can be decreased by time-shifting some of the current from a peak period to an off-peak period, even when accounting for the losses due to storage. Also, during off-peak periods, temperature and therefore resistance are typically lower, yielding additional efficiency gains [41]." \cite{lund_review_2015}


\illu{Creating a comprehensive flexibility services table by time scale etc. Basically a table version of the categories above}



\subsection{Electric Vehicles as Flexibility Resources}

"Electrical vehicles can be considered as mobile ESS." \cite{degefa_comprehensive_2021}

\subsubsection{Electric Vehicle Adoption and Charging Characteristics / Challenges}


* EV adoption trajectory and market penetration: \\
** Global and Norwegian growth trends with specific penetration statistics \\
** Future growth projections and implications for electricity demand 

* Grid impact challenges from uncontrolled EV charging: \\
** Peak load amplification and generation adequacy concerns \\
** Infrastructure constraints and potential grid violations \\
** Economic implications of increased peak demand (reduced security of supply, expensive grid expansions, reduced power quality)

* EV charging behavior patterns and availability windows: \\
** Typical charging patterns across different use cases (residential, workplace, fleet) \\
** Idle time statistics and availability for grid services \\
** Charging duration vs. connection time disparity

* Economic paradigm shift: mobile energy storage resources: \\
** Capital cost advantage over dedicated grid storage \\
** High availability (96\%) vs. low transportation utilization (4\%) \\
** Bridge to next section on smart charging solutions

\hrule

"Passenger EV sales have been growing rapidly in North America, Europe, Asia and other regions since 2012, with a cumulative two million plug-in vehicles sold by the end of 2016 (IEA, 2017b). This remains a small share of the total global vehicle stock (0.85\%), but exponential growth is occurring in many markets. It is also worth noting that the number of electric two-wheelers on the road is significantly higher than the number of EVs, with more than 200 million units already currently in circulation, mostly in China (IEA, 2016e)." \cite{international_energy_agency_status_2017}

"An increase in peak load demand could cause the need to utilize more expensive generation units, if available. Furthermore, presuming the electricity generation capacities are sufficient, the electricity network's physical constraints could be violated by the additional PEV demand" \cite{rashid_a_waraich_plug-hybrid_2013}

"The impact on the system could be positive or negative depending on a number of factors. If, for example, a large share of EV owners charge their vehicles during the evening electricity demand peak, the peak would increase, affecting generation adequacy and ramping levels." \cite{emanuele_taibi_power_2018}

"Use of the electric vehicles (EVs) is continuously increasing. EVs are high-capacity loads. Charging behavior of EVs has high level of uncertainty and produces a negative impact on PSF. However, if EVs are guided for charging in an order and combined with DSM or DR, these can be used as a flexibility improvement option with distributed energy storage devices [52]." \cite{kaushik_comprehensive_2022}.

"The impact of electric vehicles depends on the level of EV deployment and on the charging strategy employed." \cite{emanuele_taibi_power_2018}

"Most of its time, a vehicle is idle and thus offers an option for charging or discharging if equipped with a battery [305]. The charging window of EVs and PEVs is relatively long, 8–12 hours, and the charging duration relatively short, around 90 minutes, offering considerable flexibility [78]." \cite{lund_review_2015}

"Vehicles are designed to have large and frequent power fluctuations, since that is in the nature of roadway driving. The high capital cost of large generators motivates high use (average 57\% capacity factor). By contrast, personal vehicles are cheap per unit of power, and are utilized only 4\% of the time for transportation, making them potentially available the remaining 96\% of time for a secondary function." \cite{kempton_vehicle--grid_2005}


\subsubsection{Smart Charging and V2G Technologies}


* Current charging technology landscape and infrastructure state: \\
** Level 1, Level 2, DC fast charging power requirements and applications \\
** Existing infrastructure limitations and deployment status 
*** Very important to get the separation between 1. currently viable tech, 2. on the precipice of being viable, 3. future theoretical tech

* Emerging bidirectional charging standards and compatibility: \\
** These chargers both allow for bidirectional charging: CHAdeMo bidirectional capability vs. CCS evolution (CCS2 2025 timeline) \\
** Actual V2G cars: Nissan, hyunday, mitsubushi, ford, kia. Others do still not offer V2G compatible EV's, but are expected to do so in the future (Volkswagen, BMW, Volvo). \\
\research{\url{https://www.gridx.ai/knowledge/vehicle-to-grid-v2g-and-vehicle-to-home-v2h}}
** Technical requirements for V2G-enabled infrastructure \\
** Grid-forming vs. grid-following inverter implications for stability services / synthetic inertia

* Smart charging strategies and control paradigms: \\
** Controlled charging for demand-supply alignment and VRE integration \\
** Centralized vs. decentralized control architectures \\
*** NTNU is a microgrid with its own fleet of service vehicles - could use that as a case for centralized control structure \\
** Communication protocols and system integration requirements

* Economic mechanisms and market integration: \\
** Time-of-use pricing effectiveness and limitations (dual tariff paradox). Incentivize off-peak charging, but is unpredictable \\
** Grid expansion needs and associated costs \\
** Infrastructure (grid) investment deferral through covering the gap with flexibility solutions \\
** Aggregation benefits and scale requirements for market participation of EVs

* V2G technical implementation and service provision: \\
** Three core technical requirements (connection, communication, control) \\
** Power limitation factors and capability assessment \\
** Service portfolio: regulation, reserves, peak power applications

* Bridge to next section


\hrule

"Uncontrolled ("dumb") charging creates demand peaks corresponding to morning and evening traffic patterns. Time-of-use pricing (dual tariff) can actually worsen grid impacts by creating more extreme demand peaks. Smart charging strategies effectively distribute load to prevent grid constraint violations" \cite{rashid_a_waraich_plug-hybrid_2013}.

"The primary enabling function of EVs in power system transformation relates to intelligent charging, responding to when the grid has abundant generation available. Fleets of EVs can absorb excess VRE when it is available, for example through smart charging protocols. They can also serve to balance load, reducing the need for often-expensive peaking generation capacity. EVs, by serving as a new source of consumer demand, could also help overcome the difficulties that utility business models face with growing use of rooftop solar and other distributed energy resource options that affect their revenues. It is also worth noting that the widespread deployment of EVs without smart charging technology could create new challenges for system operators if low-voltage power lines become overloaded, or if charging occurs en masse during peak hours." \cite{international_energy_agency_status_2017}

"One approach which is used by many utilities today in order to shift load is a dual tariff strategy, also known as 'time of use' (TOU) pricing." \\
"Electricity consumption is much lower at night than during the day. In order to give people an incentive to shift their consumption (use of washing machines, etc.) to later times, i.e., off-peak hours, the price of electricity is low during the night" \cite{rashid_a_waraich_plug-hybrid_2013}.




"The basic concept of vehicle-to-grid power is that EDVs provide power to the grid while parked. Each vehicle must have three required elements: (1) a connection to the grid for electrical energy flow, (2) control or logical connection necessary for communication with the grid operator, and (3) controls and metering on-board the vehicle." \cite{kempton_vehicle--grid_2005}

"How much V2G power can a vehicle provide? Three independent factors limit V2G power: (1) the current-carrying capacity of the wires and other circuitry connecting the vehicle through the building to the grid, (2) the stored energy in the vehicle, divided by the time it is used, and (3) the rated maximum power of the vehicle's power electronics. The lowest of these three limits is the maximum power capability of the V2G configuration." \cite{kempton_vehicle--grid_2005}

"By establishing two-way communication between the vehicle owner and the grid, charging patterns can be manipulated to the benefit of the system" \cite{emanuele_taibi_power_2018}

"- Centralized smart charging: "One way is to give more control to the utilities that own the electricity grid. The utilities could receive information from car owners, such as where and how long their vehicles will remain parked and their expected energy consumption for the rest of the day. Based on this information, a central utility controller could determine for the PEVs when and when not to charge. In a V2G scenario, the controller could also determine when to feed power back to the energy system.
- Decentralized smart charging: publish a charging and discharging rate scale based on time and locality. Software in the PEVs could then decide when to charge or discharge electricity. An algorithm used to make such decisions would depend on data, similar to that needed in the previous case: parking duration and the location of the next activity." \cite{rashid_a_waraich_plug-hybrid_2013}


"V2G inverters operate as active power filters using self-tuning filter based control algorithms. The system simultaneously performs Grid to Vehicle charging, V2G power delivery and harmonic compensation through shunt active filtering."
\research{\url{https://www.sciencedirect.com/science/article/pii/S2352152X23030359}}


"Four-leg inverter topology, with independent neutral current control capability (the fourth leg), suppresses voltage deviations through controlled neutral point voltage manipulation while virtual impedance control schemes coordinate power sharing among multiple V2G units. 
\research{\url{https://www.sciencedirect.com/science/article/pii/S0360544222002572}}

"Inverters supporting four-quadrant operations enables both leading and lagging reactive power compensation. V2G systems function as Static Synchronous Compensators, providing dynamic reactive power injection based on real-time grid requirements."
\research{\url{https://www.researchgate.net/publication/224228227_Single-phase_inverter_design_for_V2G_reactive_power_compensation}}




"V2G could include services such as scheduled energy, power quality, reserve power, regulation, emergency load curtailment, energy balancing, and RE integration [306–309], but requires the suitable equipment (e.g. telemetry, two-way communications) and the assistance of aggregators [78,310] that lump several PEVs together to achieve the required scale. Some services may require discharging of the battery, which causes additional loss of cycle life [311] and an economic disadvantage [312]. Varying the charging power while keeping it positive allows PEVs to provide demand response and grid reliability services without causing additional wear to the battery [78]. Different smart charging schemes may be employed to minimize adverse effects from V2G to the energy system or battery [308,309,313,314]." \cite{lund_review_2015}

"A key challenge to power system service provision with ESUs is that individual units might not be available to provide the contracted service over the entire contract duration because they must also provide their local service. Therefore, there is a benefit to resource aggregation." \cite{olivier_megel_scheduling_2015}

"grid-forming converter technology has become commercially available for battery energy systems, HVDC connections and STATCOMs, but also in these applications, the grid-forming functionality is not standardized." \cite{noauthor_nordic_2023}

"The technical advancement extends to grid-forming capabilities, where inverters can create their own voltage and frequency references rather than following the existing grid. Virtual synchronous generator controls emulate the behavior of traditional rotating machines  providing synthetic inertia and enhanced stability support. This technology enables black-start capability and operation in weak grid conditions where traditional grid-following inverters would fail."
\research{\url{https://www.energy.gov/eere/solar/solar-integration-inverters-and-grid-services-basics} \&
\url{https://www.researchgate.net/publication/224408026_Static_synchronous_generators_for_distributed_generation_and_renewable_energy}}



"By utilising vehicle and charging assets when the vehicles are not in use, V2G can provide valuable services to a range of energy stakeholders, from transmission and distribution network operators to energy suppliers." \cite{aunedi_whole-system_2020}


- "In the case of vehicle-to-grid (V2G, i.e., vehicles can also feed electricity from the battery back to the grid, if adequately compensated), production cost modelling can be used to optimise V2G charging/discharging considering both the needs of the vehicle owners and system requirements for flexibility." \\
- "EVs can be modelled as batteries with constraints on availability and state of charge. For example charging/discharging is optimised so that the system is allowed to access an EV's battery as needed when it is parked (for example, at night or during working hours), but in a way that will not affect an owner's need to have the EV sufficiently charged for its mobility requirement, as this is the main purpose of an EV."\\
- "Battery EVs can provide significant flexibility when connected to the grid through smart charges, and are expected to represent the majority of battery deployment globally." \cite{emanuele_taibi_power_2018}

"The relatively lower capital costs of vehicle power systems and the low incremental costs to adapt EDVs to produce grid power suggest economic competitiveness with centralized power generation. On the other hand, compared with large generators, vehicles have low durability (about 1/50 of the design operating hours) and high cost per kWh of electric energy, suggesting that V2G power should be sold only to high-value, short-duration power markets. As we’ll describe shortly, these power markets include regulation, spinning reserves, and peak power." \cite{kempton_vehicle--grid_2005}

"EV's can fit into "fast response" and "medium response". The reason they may not be suitable for fast response are the users, as it is not operational (connected to the grid) at all times." \cite{salman_review_2022}

"There are various potential applications for V2G technology: If the demand for electricity in the grid exceeds the supply during peak hours, PEVs could supply peak power." \cite{rashid_a_waraich_plug-hybrid_2013}

"the grid operator sends requests for power to a large number of vehicles. The signal may go directly to each individual vehicle, schematically in the upper right of Fig. 1, or to the office of a fleet operator, which in turn controls vehicles in a single parking lot, schematically shown in the lower right of Fig. 1, or through a third-party aggregator of dispersed individual vehicles’ power (not shown). (The grid operator also dispatches power from traditional central-station generators using a voice telephone call or a T1line" \cite{kempton_vehicle--grid_2005}



\subsubsection{Implementation Barriers and Deployment Challenges}


\emph{Technical Barriers}
* Infrastructure compatibility and standardization challenges: \\
** OCPP 2.0 protocol implementation gaps, as well as ISO 15118/IEC 15118 \\
** IEEE1547-2018 and IEC61850 compliance requirements/standards \\
** Charging infrastructure bidirectional capability limitations \\
\research{Are current chargers capable of delivering V2G, if the car can? Are there needs for increasing transfer capacity from home/public etc. chargers?}

* Vehicle and battery constraints: \\
** Current EV fleet V2G compatibility status \\
\research{Are there any cars currently with V2G compatibility? - Which ones? Are there plans on more?}
** Battery degradation concerns and life cycle impacts \\
** State-of-charge visibility and communication protocols - privacy concerns

\emph{Regulatory and Standards Barriers}
* Definitional and standardization challenges: \\
** Terminology inconsistencies across stakeholders \\
** Grid code compliance for bidirectional power flow \\
** Safety and interconnection standards


\hrule

"Existence of incoherent terminologies and definitions is natural when new concepts and technologies are under development. Nevertheless, at certain level of technology maturity, the usage of terms and concepts amongst stakeholders will have significant implications, with impacts ranging from economic, legal and information flow." \cite{degefa_comprehensive_2021}

"On short term, variable renewable electricity production such as solar and wind with close to zero variable costs will reduce the electricity price on the wholesale spot market. Much of VRE production has also been supported by the feed-in-tariffs meaning that, all VRE production has in practice been fed to the market [364]. Even negative spot prices have been witnessed during high renewable production and low demand [364,365]. On the German market each additional GWh of VRE power fed to the grid has lowered the spot market price of electricity by \$1.4–1.7/MWh [65]. However, the consumer electricity price may rise due to the VRE feed-in-tariff: in Germany, it has increased over 20\% from 2008 to 2012 [366]." \cite{lund_review_2015}

"V2G has been studied across multiple markets [9–13], showing that EDVs cannot provide baseload power at a competitive price. This is because baseload power hits the weaknesses of EDVs—limited energy storage, short device lifetimes, and high energy costs per kWh—while not exploiting their strengths—quick response time, low standby costs, and low capital cost per kW." \cite{kempton_vehicle--grid_2005}

"Peak power is generated or purchased at times of day when high levels of power consumption are expected—for example, on hot summer afternoons. Peak power is typically generated by power plants that can be switched on for shorter periods, such as gas turbines. Since peak power is typically needed only a few hundred hours per year, it is economically sensible to draw on generators that are low in capital cost, even if each kWh generated is more expensive. Our earlier studies have shown that V2G peak power may be economic under some circumstances [9,11,14,15]. The required duration of peaking units can be 3–5h, which for V2G is possible but difficult due to on-board storage limitations. Vehicles could over come this energy-storage limit if power was drawn sequentially from a series of vehicles." \cite{kempton_vehicle--grid_2005}


"So why then have car manufacturers experienced so many problems implementing this technology on a wide scale? It is largely because standardized protocols for Battery Management Software (BMS) and Electric Vehicle Charging Stations (EVCS) have been missing. Most pilot projects have been proprietary and few have relied on ISO 15118. An interoperable language across all EV and EVCS will drastically reduce complexity and increase adoption. 

Cost also remains a barrier, as the logic managing bidirectional energy flows contradicts traditional battery management systems and still requires significant development efforts. Inverters with bidirectional capabilities also have higher hardware costs. Cost barriers can, however, be overcome in the long-run through economies of scale.

In addition, network operators are not interested in having multiple participants in the balancing energy market, which is what V2G is mostly expected to be used for. This means that they tend to either hinder the advancement of the topic or, at best, stay idle. Finally, it is not yet clear, according to regulation, who the guarantee holder would be for cars used in V2G or V2H setups. For example, who must pay if the car malfunctions while operating in a V2G mode? Such issues have to be resolved by regulators and market players before the technology can scale." \research{\url{https://www.gridx.ai/blog/how-v2g-v2h-reach-scale}}


"About the IEEE 1547-(2018) standard: Smart inverters operating under IEEE 1547-2018 standards must provide reactive power capability across their operating range, typically maintaining ±44\% of rated power as reactive capability even when producing full active power. \research{\url{https://www.esig.energy/wiki-main-page/reactive-power-capability-and-interconnection-requirements-for-pv-and-wind-plants/}}

"The IEEE 1547-2018 standard represents a paradigm shift, requiring distributed energy resources to provide enhanced grid support through mandatory voltage and frequency ride-through capabilities, reactive power capability across operating ranges, and communication protocols for coordinated control. Low voltage ride-through requirements mandate that renewable systems remain connected during voltage dips and provide reactive current injection during faults, preventing cascading disconnections that previously plagued power systems."
\research{\url{https://solarbuildermag.com/inverters/smart-pv-inverter-overview-ieee-1547-2018-and-ul-1741-explained/} \&
\url{https://www.researchgate.net/publication/266140472_Reactive_Power_Injection_Strategies_for_Single-Phase_Photovoltaic_Systems_Considering_Grid_Requirements} \&
\url{https://www.nrel.gov/grid/ieee-standard-1547/voltage-reactive-power-control}}
\research{Why is Standard IEEE 1547 not in Europe? Do they have an alternate standard?}





\subsubsection{Technical- vs. market-accessible flexibility}


* Establish the flexibility hierarchy framework: \\
** Potential -> actual -> reserves -> market-accessible progression \\
** Physical existence vs. controllability vs. economic viability vs. market participation

* Economic viability and market structure barriers for accessing flexibility: \\
\research{Keep the main focus on the Norwegian market structures (Statnett, Nord Pool). Research these structures}
** Minimum bid sizes and aggregation requirements limiting small resources \\
\research{Statnett's market participation requirements for flexibility/balancing/frequency control services}
** Uncertainty penalties of delivery (stochasticity in EV plugged in status - might infer penalties without guaranteed availability of promised flexibility). \\
** Temporal mismatches between market timeframes and technical capabilities \\
** Liquidity challenges in short-term flexibility markets (doesn't EV V2G help increase the liquidity?) \\
** Regulatory frameworks not designed for bidirectional energy storage

* Energy storage specific market integration challenges: \\
** Dual demand/supply function creating regulatory complexity \\
** Compensation mechanism inadequacies for multiple service provision ("Value-stacking", EVs being able to provide multiple flexibility solutions) \\
** Price distortions affecting storage investment economics (??)

* Aggregation as market access enabler: \\
** Aggregation of neighborhoods, parking lots, offices etc. to be able to provide enough flexibility to participate \\
** German VPP (virtual power plant) model for distributed resource integration \\

* Implications for EV flexibility deployment: \\
** Technical potential vs. current market accessibility constraints (without, and then with an aggregator) \\
\research{Are there any articles covering the tech vs. market flexibility gap? Are there any calculating methods for this?}


\hrule

"The sheer existence of flexibility resources is not sufficient in its own. Market structures and regulatory instruments are key in availing technical potential. [...] The actual tool (or combinations of tools that can be used [to increase market available flexibility] are dependent on regulatory framework in the country, the degree of decentralization in each country and the local situation." \cite{degefa_comprehensive_2021}
\skriv{expand upon the quote from Aricle 14 \cite{emanuele_taibi_power_2018} which is briefly discussed in the flexibility definitions section - about how markeds can serve as hindrances to fully realizing existing flexibility solutions}

"* Potential flexibility resources: The flexibility resources exist physically, but lacks controllability and also observability \\
* Actual flexibility resources: Flexibility exist physically and there is controllability and observability and consequently the resource is ready to be used \\
* Flexibility reserves: The part of the actual flexibility resources can be used economically [for profit/are economically viable?] \\
* Market-available flexibility reserves: The part of the flexibility reserves that can be procured from power or ancillary services market" \cite{degefa_comprehensive_2021}

"Though energy system flexibility is often perceived as a technological issue only, it has a strong link to the electricity market as well. For example, different power tariffs may enhance flexibility. On the other hand, fuel-less energy sources such as wind and solar power have a zero-marginal cost on the market, meaning that large-scale RE schemes would drop the average electricity price. Poor market design may limit access to technical flexibility in energy systems [258]." \cite{lund_review_2015}

"Integration of electricity markets, both spatially and the spot, intraday and reserve markets in a given area, could increase flexibility. Reserve market auctions should be aligned with the spot and intraday markets: if plants are committed for reserve for a much longer time than the electricity market timeframe, flexibility will be reduced [364]. Aggregation of markets over large areas provides access to more generators and loads to provide flexibility, brings about geographical smoothing, and allows for sharing of reserves, which reduces costs [258]. An auction mechanism for interconnector capacities between European power markets would increase the flexibility of the entire system [364]. However, the renewable production integration cost reduction from cross-border transmission depends on the generation mix in the neighboring countries: if VRE is largely used in both areas, the cost-saving potential of the interconnection decreases [374]. Moreover, the correlation of VRE production in the interconnected areas strongly affects the generation cost reduction due to an interconnector [375]. International market integration is being pursued in Europe with a pilot project for day-ahead market coupling recently started by 15 countries [376]." \cite{lund_review_2015}

"As the forecast errors generally decrease with a shorter prediction horizon [339,369,370], more frequent trading and reserve procurement allows for lower regulation costs and increased incomes for producers of renewable electricity [369,371], though some additional costs may also occur [371]. However, more frequent energy trading is not straightforward as intra-day markets are prone to low liquidity [367,371]. Intra-day auctions could be an attractive option for increasing intra-day market liquidity [367]." \cite{lund_review_2015}

"A major question in market design for integration of VRE is whether VRE production should follow the same market rules than the other production after it has become competitive. The same market rules could be incorporated if dynamic retail pricing schemes were in place to cover capital costs for capacity investment [373]. When the penetration of variable renewable production increases, there will be an increasing need to allocate some balancing responsibility to VRE as well, which may require special measures for balancing markets [379]. With lack of demand response and increasing VRE share, energy-only markets may not be sufficient to ensure that flexible energy plants could recover their costs for which reason some kind of capacity mechanism may be needed [373,380]. The situation could be corrected with a price-based or quantity based capacity system. In the former, the price of capacity per MW is fixed; in the latter, the amount of desired capacity is fixed, either as installed capacity or operating reserve. However, if capacity bids are in use on the reserve market, the energy price on reserve markets may get lower than on the intra-day market, creating distorting incentives [367]. As it is difficult for VRE to provide capacity, implementing a capacity mechanism would most likely lead to different market rules between VRE and conventional generation [373]." \cite{lund_review_2015}

"The EcoGrid EU project (2011-2015) worked on the development and testing of new market concept allowing to improve the balancing mechanisms by introducing a 5 minutes real-time price response to provide additional balancing power from smaller customers directly to the Transmission System Operators [53]." \cite{emil_hillberg_flexibility_2019}

"The methodology for compensating DER, in particular distributed generation, is a strong driver for uptake of these technologies. Setting the right level and structure of remuneration for grid injection is a complex undertaking with important implications. If set too high, a disproportionate amount of money will flow to DER owners; if too low, compensation might be unfair to DER owners. Fixed remuneration (per unit of energy) provides investment certainty, whereas variable pricing can more effectively encourage system-friendly VRE design choices that maximise self consumption or production during certain hours of higher system demand. Traditional compensation mechanisms, such as net energy metering, were designed on the supposition that the grid can act as a buffer for the differences in timing of electricity production and consumption of individual households. Household production and consumption are brought together on the final electricity bill. Under net energy metering, localised electricity production is implicitly valued at the rate of the variable component of the retail tariff, as the household can bank production both within and between billing periods (IEA, 2016e).25 Net billing applies a similar method, whereby injected surplus electricity is deducted from the electricity bill at a predetermined rate. In jurisdictions where a large proportion of retail tariffs consists of volumetric rates, net energy metering has come under pressure as DER owners are able to disproportionately offset their contribution to network cost. A third compensation mechanism for DERs is the feed-in tariff. In this arrangement, all electricity injected into the grid is compensated at an administratively determined rate. As pointed out in Chapter 2, many jurisdictions are shifting to other, value-based compensation for decentralised generation. Methods for such value-based compensation for DER generally fall into two categories. The first category involves taking a snapshot of current DER value, and then providing a long-term compensation guarantee based on that. A value of solar (VoS) tariff assigns fixed price tariffs based on an assessment of value components, including energy services, grid support and fuel price hedging, among others (Figure 5.7). Minnesota became the first US state to adopt a VoS tariff, with a 25-year inflation-indexed tariff that was determined through benefit-cost analysis and an extensive stakeholder consultation process (Farrell, 2014). The second category of value-based DER compensation involves more granular DER tariffs that reflect market conditions at specific times and locations. Adding price variability based on time and location can contribute to lower system costs by sending appropriate price signals to DER customers." \cite{international_energy_agency_status_2017}

"In Germany, virtual power plants are used to aggregate many DER systems and sell their cumulative excess power in real-time electricity markets. This improves the overall responsiveness of small-scale DER to market signals, as more of the available flexibility is used. At the same time, the balancing responsibility is shifted away from individual DER customers to aggregators who can better manage this risk on the basis of a portfolio of DER clients." \cite{international_energy_agency_status_2017}

"Location-dependent pricing could be used to mitigate congestion due to intermittent generation [368,373]. Nodal pricing appears to be the only option that would reflect the state of the physical power system at all times, as it is impossible to achieve that with zonal pricing. This principle applies to both domestic and international transmission capacity allocation [373]. Price caps if being high enough could facilitate fixed cost recovery by peaking units, and price floors should be low enough to reveal the value of flexibility, e.g. allowing negative prices [373]." \cite{lund_review_2015} 

"Energy storage was recognized as a high potential technology for energy system flexibility, but its role on the market is somewhat complex as it has both a demand and supply function and therefore it does not fit well into existing regulatory frameworks [381]. This in turn may discourage investments in energy storage [382–384]. Furthermore, price distortions, grid fees and the lack of price transparency may create a negative cost impact on energy storage [9]. IEA suggests that policymakers should enable compensation for the multitude of services energy storage provides, e.g. payment based on the value of reliability, power quality, energy security and efficiency gains. Potential mechanisms include real-time pricing, pricing by service and taxation being applied to final products. In the U.S., efforts have been made to enable incentives to energy storage, e.g. through opening electricity markets to energy storage and permitting companies other than large utilities to sell ancillary services [9]." \cite{lund_review_2015}

"Price-based DSM increases demand elasticity and hence it could reduce price spikes [31]. DSM can also shift market power from generators to consumers [37], reduce prices [385] and price volatility [50], and make the market more efficient [385]. In many electricity markets where VRE has priority, price is negatively correlated with VRE production [35] and hence price-based DSM can provide flexibility for VRE integration directly. However, the closed-loop feedback system between the physical and market layer created by real-time pricing, together with increasing VRE production, may lead to increased volatility [386]. Large-scale system models including the interactions between the market and physical system are needed for full understanding of DSM effects [37]. Market rules affect the possibility of DSM to enter different markets. For example, auctions on primary and secondary reserves in Germany are held 1 month and 6 months before delivery, which may be too restrictive for DSM providers [33]." \cite{lund_review_2015}


\subsection{Current studies/pilot projects on Microgrids and Local Energy Systems}

* Small, local (micro) grid EV V2G/smart charging flexibility studies: \\
** Methodological precedents for campus-scale EV integration modeling \\
** Performance metrics used in similar distribution grid studies \\
** Modeling tools and validation approaches from comparable research \\
** V2G fleet studies: Quantifies operational cost reductions and grid benefits \\


* Real-world V2G pilot projects with quantified benefits: \\
** Living Lab Smart Charging - deployment scale and measured impacts: \\
\research{Equans/GreenFlux project \url{https://www.greenflux.com/expertise/blogs/equans-innovative-living-lab-at-dordrecht/}}

** Norwegian distribution grid flexibility cases (ferry charging): \\
\research{Electric ferry charging in Norway, from Article 22. Charging need surpassing current trafo capabilities - Are there any studies/publications on this?}

** NODES project: \\
\research{NODES proof of concept on: "A new marketplace capable of exploiting decentralised flexibility"} 

** CURRENT has a Smart Charging Alignment for Europe (SCAlE) project
\research{Read about the SCAlE project}

** Eurostar has a project regarding movable EV chargers with V2X compatability (SpotMe)
\research{Read about the SpotMe project}

\research{Read about the FuChar project}
\research{Are CINELDI having projects about this?}

"The Edison Project intends to install enough turbines to accommodate 50\% of Denmark's total power needs, while using V2G to protect the grid. The Edison Project plans to use PEVs while they are plugged into the grid to store additional wind energy that the grid cannot handle. During the hours of peak energy use, or when the wind is calm, the power stored in these PEVs will be fed into the grid. To aid in the acceptance of PEVs, zero emission vehicles received government subsidies.[citation needed]

Following the Edison project, the Nikola project was started[54] which focused on demonstrating V2G technology in a laboratory setting at the Risø Campus of the Technical University of Denmark (DTU). DTU is a partner along with Nuvve and Nissan. The Nikola project was completed in 2016, laying the groundwork for the Parker project, which used a fleet of EVs to demonstrate the technology in a real-life setting. This project was partnered by DTU,[55] Insero, Nuvve, Nissan and Frederiksberg Forsyning (a Danish distribution system operator in Copenhagen). The partners explored commercial opportunities by systematically testing and demonstrating V2G services across car brands. Economic and regulatory barriers were identified as well as the economic and technical impacts of the applications on the power system and markets.[56] The project started in August 2016 and ended in September 2018." \\
\research{\url{https://en.wikipedia.org/wiki/Vehicle-to-grid}}
\research{The Edison project, The Nikola project and the Parker project in Denmark. Wind supported by V2G EVs.}


* Research gap identification and thesis positioning: \\
** Limitations in existing campus-scale studies \\
** data quality and modeling assumtion gaps (in current studies?) \\
** This thesis contribution: Utilizing actual NTNU infrastructure and charging data for realistic modeling 

\nb{All articles included here should support one of three conditions: 1 - my choice of model (structure, tools, assumptions). 2 - My choice of parameters, which to include and which to not. 3 - For comparing and/or validation purposes regaring my results}

\hrule

"A new marketplace capable of exploiting decentralised flexibility is being developed by the Norwegian energy company Agder Energi and the Power Exchange Nord Pool through the recently established subsidiary NODES [63]. The intention is that this new marketplace, placing a value on flexibility, will bridge the gap between current wholesale power markets and local flexibility markets. The proof of concept has been demonstrated, and the intention is to utilise entities at different levels of the power system to provide and utilise flexibility services through a transparent marketplace open to all potential market participants, illustrated in Figure 7. Using available flexibility is an alternative to grid investments and creates opportunities for customers providing new services and enables greater integration of renewable energy." \cite{emil_hillberg_flexibility_2019}

"In the Netherlands, the “Living Lab Smart Charging” project is a large-scale open platform to support adoption of smart charging and vehicle-to-grid technologies. EV owners can earn money by allowing their vehicles to be charged according to a combination of their needs and what the grid needs (Lambert, 2016)." \cite{international_energy_agency_status_2017}

"Operating cost reduces substantially as V2G enabled fleets displace up to 50\% of FR otherwise provided by part-loaded and therefore less efficient thermal generators. This also results in lower renewable curtailment due to fewer FR-providing" \cite{aunedi_whole-system_2020}

"a simple use case relevant to Norwegian DSOs: Flexibility resources as a measure to support the integration of electrified maritime transportation. Infrastructure for charging of electrical ferries is being installed in several small Norwegian coastal towns or villages that are supplied by distribution grids with insufficient power capacity for the power demand peaks during charging. See Fig. 7 for an illustration. As an example, the area may have a base load demand around 2 MW, but charging ferry when at quay (for approximately 7 min) requires an additional 4 MW. If the grid capacity is 5 MW, there is either a need for congestion management services or for costly grid reinforcement measures." \cite{degefa_comprehensive_2021}
\illu{From: Article 22, Figure 7 (page 13) - \cite{degefa_comprehensive_2021} . Good simple illustration about the electric ferry case}



\section{Methodology}
\subsection{Case Study: NTNU Gløshaugen Campus}

\subsubsection{Overview of NTNU's Electrical infrastructure}

* Detailed description of the campus electrical network, including: \\
** Network topology(assume all 24 kV(?) cables, since the school is going to change them all in the future. \\ 
** Distribution system characteristics. \\
** Key infrastructure components (transformers, cables etc). \\
** Building energy usage, how close to capacity etc. 


\subsubsection{Existing EV Charging Infrastructure}

* Information on current EV charging facilities at Gløshaugen: \\
** Number and types of chargers. \\
** Locations and distribution across campus. \\
** Technical specifications(power ratings, connectivity)


\subsubsection{Data Collection}

* Sources of data used in the study: \\
** Electrical load profiles for the campus (driftsentralen). \\
** Historical EV charging data (Current.eco). \\
** Grid topology (driftsentralen). \\

* Discussion of data processing methods and assumptions made


\subsection{Simulation Tools and Models}

* PowerFactory + Python dataanalysis (justify why these are used and good to use)

* Specify library versions (specially Python)

* Load flow analysis, Time-series simulations, integration with EV charging models


\subsection{Modeling Assumption, Parameters and Uncertainty}

* Outline key assumptions, f.ex.: \\
** EV arrival and departure times(different for private cars and service vehicles). \\
** Charging/discharging rates and battery capacities. \\
** User participation rates, changes in different scenarios \\
** The conductor changes reactance based on weather conditions (hva faen er det? Grid-ledningen??). Are they in the ground? Are they in the air? Temperature, more?? \\

* EV and charging needs assumtions: \\
\research{How much power do the different cars need (Hiace, truck, employee)}
** Employee with their personal car needs the battery to be on a certain percentage at the end of the day (lets say minimum 70\%) \\
** Service cars owned by NTNU needs to be at full charge at the beginning of the working day \\
** Assume that all employees have the same car, and that NTNU has a hiace type car and some trucks(?)

* Model EVs as controllable reactive power sources (±Q capability) \\
* Four-quadrant operation enables voltage regulation during all charging states
\research{\url{https://ieeexplore.ieee.org/document/8073954}}



* Justification of assumptions and parameters based on data. 

* Consider including uncertainties. Quantify these and provide justification as to why they are simplified away, or not included as a (stochastic) modeling parameter (weather/temperature, time of arrival, SOC after use). Rather, we opt for doing a sensitivity-analysis on these uncertain parameters.

\subsection{Performance Metrics (for the Results)}

* Metrics for assessing the impact of EV integration on the network. \\ 
** Peak Load \\
** Voltage stability indices \\
** Power quality indices \\
** Energy losses \\

Egne forslag: \\
** Voltage profiles and deviations. \\ 
** Transformer and feeder loading levels. \\
** Energy losses and efficiency. \\
** Active-Reactive power space???. \\
** SAIDI / SAIFI / CAIDI reliability indices. \\
** Money saved \\ 

Claude: \\
* Voltage deviation $\pm$ \% from nominal \\
* Transformer loading/overloading - \% of rated power capacity \\
* Economic metrics - money saved/earned by EV users


\subsection{Simulation Scenarios}

* Consider looking into scenarios where demand increases drastically (another building is attached to the campus grid [justified, as Gløshaugen is expanding due to the Campus Project]). And then compare how the grid is capable of handling these increases with and without V2G support from its charging points. - Is it enough to deffer some grid expansion on the campus? 

* Consider adding some stochasticity in the demand profile, to see how the V2G capabilities handle those, compared to the system as it stands now

\subsubsection{Base case Scenario}

* MAYBE: Model the campus without the energy consumption used on EV charging

* Model the campus network as is at current date, with uncontrolled charging


\subsubsection{High Utilization of Charging Points / High EV Penetration Scenario}

* Simulating increased EV adoption rates (if possible; base of projections)

* Assessing the impact of uncontrolled charging on the network


\subsubsection{Smart charging Scenario: Load-shifting}

* Implementing controlled charging strategies in the model

* Flattening EV charging load peaks

* Evaluating effectiveness in mitigating impacts of increased EV penetration (How much does this bring the added strain on the grid back to the no-EV OR unplanned charging scenario?)


\subsubsection{V2G Capability Scenario}

* Modeling bidirectional energy flow with EVs providing services back to the grid

* Analyzing the potential benefits for network performance and flexibility

* Evaluating effectiveness in mitigating impacts of increased EV penetration (How much does this bring the added strain on the grid back to the no-EV OR unplanned charging scenario?)

* Compensation models for participants? Portion of money saved/earned by the charging company? Compensating for strain on battery ["kjøregodtgjørelse"]


\subsection{Evaluation methods for results}

* Methods for calculating and interpreting the performance metrics

* Sensitivity analysis: \\
** How variation in key parameters affect results. \\ 
** Time connected to charger, state of charge when the car is replugged into the charger after use, number of cars connected. \\
** Approach to testing the robustness of findings \\
** Monte Carlo method?



\section{Results}

\subsection{Simulation outcomes}

\subsubsection{Base Case and/or with current EV charging need}

* Presentation of network performance metrics

* Identification of existing issues/bottlenecks


\subsubsection{High Utilization of Charging Points}

* Analysis of the impact of increased EV adoption

* Identification of areas of potential issues


\subsubsection{Smart Charging Scenario Results}

* Evalutation of how it mitigates negative impacts of high EV utilization and network performance

* Comparison??

\subsubsection{V2G Scenario Results}

* Assessment of the benefits and/or challenges of this implementation

* Evaluation of changes in network performance and flexibility


\subsection{Comparative Analysis}

* Side by side comparison of all scenarios

* Visual representations to illustrate differences

* Highlight key findings and differences between scenarios


\subsection{Cost-Benefit Analysis of Proposed Flexibility Solutions}

* Estimate costs of implementing smart charging and V2G

* Quantify potential benefits (possibly: Deferred grid upgrades, reduced VRE curtailment, earnings for charging company and EV owners)

* Assess economic viability of proposed solutions


\section{Discussion}

\subsection{Interpretation of Results}

* In depth analysis of what the results imply for the Gløshaugen campus network

* Examine the effectiveness of different mitigation strategies



\subsection{Challenges and Barriers to Implementing Flexibility Solutions}

* Technical limitations, infrastructure changes needed to make this work

* Regulatory hurdles, policy changes needed to make this work

* Consumer adoption, adaptation and acceptance


\subsection{Technical vs. Market flexibility gap}

* While the technical flexibility is shown to be X (from results section), market-available flexibility is likely lower due to a,b,c policies (as discussed in theory section). [Regne ut eller ikke]

* Overgangstekst for "hvordan å fikse dette" i neste seksjon.


\subsection{Policy, Market and Infrastructure Implications}

* Consideration of how these results inform broader effects of higher EV penetration and V2G-technologies to increase the robustness of the grid.

* Discuss the possible benefits (or feasibility) of placing charging stations strategically in areas with lacking grid capacity.

* Discuss whether it can be inferred that proper utilization of V2G flexibility can defer/postpone/make some grid expansion developments obsolete because it provides enough stability/peak power.

* Discuss potential policy or infrastructure changes needed and their likelihood. 


\subsection{Scalability and Generalization of Results}

* Discussion of how the findings can be generalized to other parts of the grid, and how this affects the future of EVs, grid expansion and potential of EVs offering flexibility. 


\subsection{Limitations and Future Work}

* Acknowledgment of any limitations in the methodology or data

* How might these limitations have influenced the results

* Identify areas for improvement in modeling and/or data collections

* Suggestions for further research to build on your findings

\hrule
Expansion of model: Including current/possible future market constraints for minimum bids, penalties, aggregators


\section{Conclusion}

* Summarize the key findings and their significance

* Reflect on how the research objectives were met

* Reiterate the contributions from this thesis to the power system flexibility and EV integration


Article 1: \cite{shinde_literature_2019} \\
Article 2: \cite{salman_review_2022} \\
Article 3: \cite{lechl_review_2023} \\
Article 4: \cite{michael_emmanuel_review_2020} \\
Article 5: \cite{eamonn_lannoye_assessing_2015} \\
Article 6: \cite{impram_challenges_2020} \\
Article 7: \cite{mohseni_demand_2022} \\
Article 8: \cite{eamonn_lannoye_evaluation_2012} \\
Article 9: \cite{emil_hillberg_flexibility_2019} \\
Article 10: \cite{kaushik_comprehensive_2022} \\
Article 11: \cite{olivier_megel_scheduling_2015} \\
Article 12: \cite{noauthor_nordic_2023} \\
Article 13: \cite{rashid_a_waraich_plug-hybrid_2013} \\
Article 14: \cite{emanuele_taibi_power_2018} \\
Article 15: \cite{g_papaefthymiou_power_2018} \\
Article 16: \cite{alireza_akrami_power_2019} \\
Article 17: \cite{lund_review_2015} \\
Article 18: \cite{international_energy_agency_status_2017} \\
Article 19: \cite{kempton_vehicle--grid_2005} \\
Article 20: \cite{aunedi_whole-system_2020} \\
Article 21: \cite{noauthor_veileder_nodate} \\
Article 22: \cite{degefa_comprehensive_2021} \\
Article 23: \cite{li_grid-side_2018} 


